% Options for packages loaded elsewhere
\PassOptionsToPackage{unicode}{hyperref}
\PassOptionsToPackage{hyphens}{url}
%
\documentclass[
]{article}
\usepackage{amsmath,amssymb}
\usepackage{iftex}
\ifPDFTeX
  \usepackage[T1]{fontenc}
  \usepackage[utf8]{inputenc}
  \usepackage{textcomp} % provide euro and other symbols
\else % if luatex or xetex
  \usepackage{unicode-math} % this also loads fontspec
  \defaultfontfeatures{Scale=MatchLowercase}
  \defaultfontfeatures[\rmfamily]{Ligatures=TeX,Scale=1}
\fi
\usepackage{lmodern}
\ifPDFTeX\else
  % xetex/luatex font selection
\fi
% Use upquote if available, for straight quotes in verbatim environments
\IfFileExists{upquote.sty}{\usepackage{upquote}}{}
\IfFileExists{microtype.sty}{% use microtype if available
  \usepackage[]{microtype}
  \UseMicrotypeSet[protrusion]{basicmath} % disable protrusion for tt fonts
}{}
\makeatletter
\@ifundefined{KOMAClassName}{% if non-KOMA class
  \IfFileExists{parskip.sty}{%
    \usepackage{parskip}
  }{% else
    \setlength{\parindent}{0pt}
    \setlength{\parskip}{6pt plus 2pt minus 1pt}}
}{% if KOMA class
  \KOMAoptions{parskip=half}}
\makeatother
\usepackage{xcolor}
\usepackage[margin=1in]{geometry}
\usepackage{color}
\usepackage{fancyvrb}
\newcommand{\VerbBar}{|}
\newcommand{\VERB}{\Verb[commandchars=\\\{\}]}
\DefineVerbatimEnvironment{Highlighting}{Verbatim}{commandchars=\\\{\}}
% Add ',fontsize=\small' for more characters per line
\usepackage{framed}
\definecolor{shadecolor}{RGB}{248,248,248}
\newenvironment{Shaded}{\begin{snugshade}}{\end{snugshade}}
\newcommand{\AlertTok}[1]{\textcolor[rgb]{0.94,0.16,0.16}{#1}}
\newcommand{\AnnotationTok}[1]{\textcolor[rgb]{0.56,0.35,0.01}{\textbf{\textit{#1}}}}
\newcommand{\AttributeTok}[1]{\textcolor[rgb]{0.13,0.29,0.53}{#1}}
\newcommand{\BaseNTok}[1]{\textcolor[rgb]{0.00,0.00,0.81}{#1}}
\newcommand{\BuiltInTok}[1]{#1}
\newcommand{\CharTok}[1]{\textcolor[rgb]{0.31,0.60,0.02}{#1}}
\newcommand{\CommentTok}[1]{\textcolor[rgb]{0.56,0.35,0.01}{\textit{#1}}}
\newcommand{\CommentVarTok}[1]{\textcolor[rgb]{0.56,0.35,0.01}{\textbf{\textit{#1}}}}
\newcommand{\ConstantTok}[1]{\textcolor[rgb]{0.56,0.35,0.01}{#1}}
\newcommand{\ControlFlowTok}[1]{\textcolor[rgb]{0.13,0.29,0.53}{\textbf{#1}}}
\newcommand{\DataTypeTok}[1]{\textcolor[rgb]{0.13,0.29,0.53}{#1}}
\newcommand{\DecValTok}[1]{\textcolor[rgb]{0.00,0.00,0.81}{#1}}
\newcommand{\DocumentationTok}[1]{\textcolor[rgb]{0.56,0.35,0.01}{\textbf{\textit{#1}}}}
\newcommand{\ErrorTok}[1]{\textcolor[rgb]{0.64,0.00,0.00}{\textbf{#1}}}
\newcommand{\ExtensionTok}[1]{#1}
\newcommand{\FloatTok}[1]{\textcolor[rgb]{0.00,0.00,0.81}{#1}}
\newcommand{\FunctionTok}[1]{\textcolor[rgb]{0.13,0.29,0.53}{\textbf{#1}}}
\newcommand{\ImportTok}[1]{#1}
\newcommand{\InformationTok}[1]{\textcolor[rgb]{0.56,0.35,0.01}{\textbf{\textit{#1}}}}
\newcommand{\KeywordTok}[1]{\textcolor[rgb]{0.13,0.29,0.53}{\textbf{#1}}}
\newcommand{\NormalTok}[1]{#1}
\newcommand{\OperatorTok}[1]{\textcolor[rgb]{0.81,0.36,0.00}{\textbf{#1}}}
\newcommand{\OtherTok}[1]{\textcolor[rgb]{0.56,0.35,0.01}{#1}}
\newcommand{\PreprocessorTok}[1]{\textcolor[rgb]{0.56,0.35,0.01}{\textit{#1}}}
\newcommand{\RegionMarkerTok}[1]{#1}
\newcommand{\SpecialCharTok}[1]{\textcolor[rgb]{0.81,0.36,0.00}{\textbf{#1}}}
\newcommand{\SpecialStringTok}[1]{\textcolor[rgb]{0.31,0.60,0.02}{#1}}
\newcommand{\StringTok}[1]{\textcolor[rgb]{0.31,0.60,0.02}{#1}}
\newcommand{\VariableTok}[1]{\textcolor[rgb]{0.00,0.00,0.00}{#1}}
\newcommand{\VerbatimStringTok}[1]{\textcolor[rgb]{0.31,0.60,0.02}{#1}}
\newcommand{\WarningTok}[1]{\textcolor[rgb]{0.56,0.35,0.01}{\textbf{\textit{#1}}}}
\usepackage{longtable,booktabs,array}
\usepackage{calc} % for calculating minipage widths
% Correct order of tables after \paragraph or \subparagraph
\usepackage{etoolbox}
\makeatletter
\patchcmd\longtable{\par}{\if@noskipsec\mbox{}\fi\par}{}{}
\makeatother
% Allow footnotes in longtable head/foot
\IfFileExists{footnotehyper.sty}{\usepackage{footnotehyper}}{\usepackage{footnote}}
\makesavenoteenv{longtable}
\usepackage{graphicx}
\makeatletter
\def\maxwidth{\ifdim\Gin@nat@width>\linewidth\linewidth\else\Gin@nat@width\fi}
\def\maxheight{\ifdim\Gin@nat@height>\textheight\textheight\else\Gin@nat@height\fi}
\makeatother
% Scale images if necessary, so that they will not overflow the page
% margins by default, and it is still possible to overwrite the defaults
% using explicit options in \includegraphics[width, height, ...]{}
\setkeys{Gin}{width=\maxwidth,height=\maxheight,keepaspectratio}
% Set default figure placement to htbp
\makeatletter
\def\fps@figure{htbp}
\makeatother
\setlength{\emergencystretch}{3em} % prevent overfull lines
\providecommand{\tightlist}{%
  \setlength{\itemsep}{0pt}\setlength{\parskip}{0pt}}
\setcounter{secnumdepth}{5}
\usepackage{booktabs}
\usepackage{longtable}
\usepackage{array}
\usepackage{multirow}
\usepackage{wrapfig}
\usepackage{float}
\usepackage{colortbl}
\usepackage{pdflscape}
\usepackage{tabu}
\usepackage{threeparttable}
\usepackage{threeparttablex}
\usepackage[normalem]{ulem}
\usepackage{makecell}
\usepackage{xcolor}
\ifLuaTeX
  \usepackage{selnolig}  % disable illegal ligatures
\fi
\IfFileExists{bookmark.sty}{\usepackage{bookmark}}{\usepackage{hyperref}}
\IfFileExists{xurl.sty}{\usepackage{xurl}}{} % add URL line breaks if available
\urlstyle{same}
\hypersetup{
  pdftitle={Análisis de Supervivencia en Pacientes que han Sufrido un Ictus},
  hidelinks,
  pdfcreator={LaTeX via pandoc}}

\title{Análisis de Supervivencia en Pacientes que han Sufrido un Ictus}
\usepackage{etoolbox}
\makeatletter
\providecommand{\subtitle}[1]{% add subtitle to \maketitle
  \apptocmd{\@title}{\par {\large #1 \par}}{}{}
}
\makeatother
\subtitle{Trabajo de Evaluación Continuada --- Segunda Parte}
\author{}
\date{\vspace{-2.5em}Curso 2024-25}

\begin{document}
\maketitle

{
\setcounter{tocdepth}{2}
\tableofcontents
}
\begin{center}\rule{0.5\linewidth}{0.5pt}\end{center}

\hypertarget{introducciuxf3n-y-contexto-del-estudio}{%
\section{Introducción y contexto del
estudio}\label{introducciuxf3n-y-contexto-del-estudio}}

El ictus es una de las principales causas de muerte y discapacidad en el
mundo. Se produce cuando el flujo de sangre hacia una parte del cerebro
se interrumpe de forma súbita, provocando un daño que puede resultar
fatal. Según la Organización Mundial de la Salud, el ictus causa
aproximadamente 5,5 millones de muertes al año a nivel global, siendo la
segunda causa de muerte más frecuente tras las enfermedades coronarias.

En este trabajo se analizan los datos de un \textbf{estudio de
seguimiento prospectivo} llevado a cabo en una red de hospitales
públicos españoles. El estudio tiene como objetivo evaluar la
\textbf{supervivencia de pacientes tras haber sufrido un ictus}, así
como identificar qué factores influyen en el riesgo de fallecimiento
durante los dos años posteriores al episodio.

Un aspecto central del estudio es la comparación entre el
\textbf{tratamiento estándar} actualmente empleado en la práctica
clínica y un \textbf{nuevo tratamiento que se encuentra en fase de
evaluación}, y que ha mostrado resultados prometedores en estudios
previos. Junto con el tipo de tratamiento, se recogen también
características básicas del paciente en el momento del ingreso que
podrían condicionar su evolución.

El evento de interés es el \textbf{fallecimiento del paciente} durante
el período de seguimiento. Los pacientes que siguen vivos al finalizar
los dos años de seguimiento (730 días) se consideran
\textbf{observaciones censuradas}, al igual que aquellos que abandonan
el estudio antes de su conclusión.

\begin{center}\rule{0.5\linewidth}{0.5pt}\end{center}

\hypertarget{descripciuxf3n-de-las-variables}{%
\section{Descripción de las
variables}\label{descripciuxf3n-de-las-variables}}

La muestra está formada por \textbf{1.400 pacientes} ingresados tras
sufrir un ictus. Para cada uno de ellos se dispone de las siguientes
variables:

\begin{longtable}[]{@{}
  >{\raggedright\arraybackslash}p{(\columnwidth - 4\tabcolsep) * \real{0.3448}}
  >{\raggedright\arraybackslash}p{(\columnwidth - 4\tabcolsep) * \real{0.4483}}
  >{\raggedright\arraybackslash}p{(\columnwidth - 4\tabcolsep) * \real{0.2069}}@{}}
\toprule\noalign{}
\begin{minipage}[b]{\linewidth}\raggedright
Variable
\end{minipage} & \begin{minipage}[b]{\linewidth}\raggedright
Descripción
\end{minipage} & \begin{minipage}[b]{\linewidth}\raggedright
Tipo
\end{minipage} \\
\midrule\noalign{}
\endhead
\bottomrule\noalign{}
\endlastfoot
\texttt{Survt} & Tiempo de supervivencia desde el ictus hasta el
fallecimiento o hasta la censura, medido en días & Cuantitativa
continua \\
\texttt{Status} & Variable de censura: \texttt{1} si el paciente fallece
durante el seguimiento; \texttt{0} si no se produce el evento &
Binaria \\
\texttt{x1} & Tipo de tratamiento: \texttt{1} = tratamiento estándar;
\texttt{0} = nuevo tratamiento en estudio & Binaria \\
\texttt{x2} & Edad del paciente en el momento del ictus, en años &
Cuantitativa continua \\
\texttt{x3} & Presión arterial sistólica registrada al ingreso
hospitalario, en mmHg & Cuantitativa continua \\
\texttt{x4} & Sexo del paciente: \texttt{1} = mujer; \texttt{0} = hombre
& Binaria \\
\texttt{x5} & Índice de masa corporal (IMC) del paciente al ingreso, en
kg/m² & Cuantitativa continua \\
\end{longtable}

A continuación se describe el papel esperado de cada variable en el
análisis:

\begin{itemize}
\item
  \textbf{\texttt{x1} --- Tratamiento:} Es la variable de interés
  principal del estudio. Se compara si los pacientes que reciben el
  nuevo tratamiento en evaluación presentan una supervivencia mayor que
  los que reciben el tratamiento estándar.
\item
  \textbf{\texttt{x2} --- Edad:} A mayor edad, el organismo tiene menor
  capacidad de recuperación y mayor vulnerabilidad ante las
  complicaciones, por lo que se espera que la edad esté asociada a un
  mayor riesgo de fallecimiento.
\item
  \textbf{\texttt{x3} --- Presión arterial sistólica al ingreso:}
  Refleja el estado cardiovascular del paciente en el momento del
  episodio. Tanto valores muy bajos como muy elevados pueden indicar
  mayor gravedad y asociarse a un peor pronóstico.
\item
  \textbf{\texttt{x4} --- Sexo:} Existen diferencias entre hombres y
  mujeres en la incidencia y evolución del ictus. Se incluye como
  posible factor de confusión o modificador del efecto del tratamiento.
\item
  \textbf{\texttt{x5} --- IMC:} El índice de masa corporal es un
  indicador del estado nutricional general del paciente. Valores
  extremos pueden estar asociados a un peor pronóstico durante el
  seguimiento.
\end{itemize}

\begin{center}\rule{0.5\linewidth}{0.5pt}\end{center}

\hypertarget{objetivos-del-anuxe1lisis}{%
\section{Objetivos del análisis}\label{objetivos-del-anuxe1lisis}}

El análisis se estructura en torno a los siguientes objetivos:

\begin{enumerate}
\def\labelenumi{\arabic{enumi}.}
\item
  Describir la supervivencia global de los pacientes mediante el
  estimador de Kaplan-Meier y comparar las curvas de supervivencia según
  el tipo de tratamiento recibido.
\item
  Ajustar un \textbf{modelo de Cox estándar} que permita estimar el
  efecto de cada covariable sobre el riesgo de fallecimiento,
  controlando simultáneamente por el resto de factores.
\item
  Verificar el \textbf{supuesto de proporcionalidad de riesgos} del
  modelo de Cox para cada covariable, mediante el test de Schoenfeld y
  el análisis gráfico correspondiente.
\item
  En caso de que alguna covariable no cumpla el supuesto de
  proporcionalidad, ajustar un \textbf{modelo de Cox estratificado} que
  permita que la función de riesgo basal varíe entre los grupos
  definidos por dicha variable.
\item
  Considerar un \textbf{modelo de Cox ampliado} que incorpore
  covariables dependientes del tiempo, como alternativa o complemento al
  modelo estratificado.
\item
  Interpretar todos los resultados en términos de \emph{hazard ratios}
  (HR) e intervalos de confianza al 95\%, en el contexto del estudio de
  supervivencia tras ictus.
\end{enumerate}

\begin{center}\rule{0.5\linewidth}{0.5pt}\end{center}

\hypertarget{planteamiento-metodoluxf3gico}{%
\section{Planteamiento
metodológico}\label{planteamiento-metodoluxf3gico}}

\hypertarget{el-modelo-de-cox}{%
\subsection{El modelo de Cox}\label{el-modelo-de-cox}}

El modelo de regresión de Cox (Cox, 1972) es el modelo semiparamétrico
de referencia en el análisis de supervivencia. Permite estudiar el
efecto de un conjunto de covariables sobre el riesgo instantáneo de que
ocurra el evento de interés (en este caso, el fallecimiento), sin
necesidad de asumir ninguna forma paramétrica concreta para dicho riesgo
a lo largo del tiempo.

Para el individuo \(i\), el modelo expresa la función de riesgo como:

\[h_i(t) = h_0(t) \cdot \exp\left(\beta_1 x_{1i} + \beta_2 x_{2i} + \beta_3 x_{3i} + \beta_4 x_{4i} + \beta_5 x_{5i}\right)\]

donde \(h_0(t)\) es la función de riesgo basal, común a todos los
individuos, y los coeficientes \(\beta_j\) miden el efecto de cada
covariable. El exponencial de cada coeficiente, \(\exp(\beta_j)\), se
interpreta como el \textbf{hazard ratio (HR)}: cuántas veces es mayor el
riesgo de un grupo respecto al de referencia, manteniendo fijas el resto
de variables.

El modelo de Cox estándar asume que los riesgos de los distintos grupos
son \textbf{proporcionales} a lo largo del tiempo, es decir, que el HR
entre dos individuos cualesquiera es constante en el tiempo. Este
supuesto debe verificarse.

\hypertarget{modelo-estratificado}{%
\subsection{Modelo estratificado}\label{modelo-estratificado}}

Si una covariable viola el supuesto de proporcionalidad, una solución es
\textbf{estratificar} el modelo por dicha variable. En este caso, se
permite que cada estrato tenga su propia función de riesgo basal
\(h_{0g}(t)\), mientras que el efecto del resto de covariables sigue
siendo común a todos los estratos:

\[h_{ig}(t) = h_{0g}(t) \cdot \exp\left(\sum_j \beta_j x_{ij}\right)\]

La variable de estratificación no genera un coeficiente propio (no se
obtiene un HR para ella), pero su efecto sobre el riesgo basal queda
controlado.

\hypertarget{modelo-ampliado}{%
\subsection{Modelo ampliado}\label{modelo-ampliado}}

El modelo de Cox ampliado permite incorporar \textbf{covariables
dependientes del tiempo}, cuyo valor o efecto puede cambiar a lo largo
del seguimiento. Una forma habitual de abordar la no proporcionalidad es
incluir un término de interacción entre la covariable problemática y una
función del tiempo, por ejemplo \(x_j \cdot \log(t)\):

\[h_i(t) = h_0(t) \cdot \exp\left(\beta_j x_{ji} + \gamma_j \left[x_{ji} \cdot \log(t)\right] + \ldots\right)\]

Si el coeficiente \(\gamma_j\) resulta significativo, se confirma que el
efecto de la covariable varía con el tiempo, explicando así la violación
del supuesto de proporcionalidad.

\begin{center}\rule{0.5\linewidth}{0.5pt}\end{center}

\hypertarget{exploraciuxf3n-inicial-de-los-datos}{%
\section{Exploración inicial de los
datos}\label{exploraciuxf3n-inicial-de-los-datos}}

\hypertarget{carga-de-libreruxedas-y-datos}{%
\subsection{Carga de librerías y
datos}\label{carga-de-libreruxedas-y-datos}}

\begin{Shaded}
\begin{Highlighting}[]
\FunctionTok{library}\NormalTok{(survival)}
\FunctionTok{library}\NormalTok{(survminer)}
\FunctionTok{library}\NormalTok{(ggplot2)}
\FunctionTok{library}\NormalTok{(dplyr)}
\FunctionTok{library}\NormalTok{(knitr)}
\FunctionTok{library}\NormalTok{(kableExtra)}
\end{Highlighting}
\end{Shaded}

\begin{Shaded}
\begin{Highlighting}[]
\CommentTok{\# Lectura del fichero de datos}
\NormalTok{datos }\OtherTok{\textless{}{-}} \FunctionTok{read.table}\NormalTok{(}\StringTok{"datos.txt"}\NormalTok{, }\AttributeTok{header =} \ConstantTok{TRUE}\NormalTok{, }\AttributeTok{sep =} \StringTok{""}\NormalTok{)}
\end{Highlighting}
\end{Shaded}

\hypertarget{resumen-descriptivo}{%
\subsection{Resumen descriptivo}\label{resumen-descriptivo}}

\begin{Shaded}
\begin{Highlighting}[]
\FunctionTok{summary}\NormalTok{(datos)}
\end{Highlighting}
\end{Shaded}

\begin{verbatim}
##      Survt             Status             x1              x2       
##  Min.   :  1.028   Min.   :0.0000   Min.   :0.000   Min.   :1.000  
##  1st Qu.:191.963   1st Qu.:1.0000   1st Qu.:0.000   1st Qu.:3.274  
##  Median :370.068   Median :1.0000   Median :0.000   Median :5.500  
##  Mean   :367.305   Mean   :0.7975   Mean   :0.433   Mean   :5.494  
##  3rd Qu.:538.874   3rd Qu.:1.0000   3rd Qu.:1.000   3rd Qu.:7.698  
##  Max.   :729.893   Max.   :1.0000   Max.   :1.000   Max.   :9.995  
##        x3                 x4              x5       
##  Min.   :-12.3540   Min.   :1.000   Min.   :1.009  
##  1st Qu.:  0.4667   1st Qu.:1.000   1st Qu.:3.197  
##  Median :  3.9620   Median :2.000   Median :5.564  
##  Mean   :  3.9547   Mean   :1.514   Mean   :5.509  
##  3rd Qu.:  7.4158   3rd Qu.:2.000   3rd Qu.:7.815  
##  Max.   : 23.4707   Max.   :2.000   Max.   :9.998
\end{verbatim}

\begin{Shaded}
\begin{Highlighting}[]
\FunctionTok{data.frame}\NormalTok{(}
  \AttributeTok{Variable =} \FunctionTok{c}\NormalTok{(}
    \StringTok{"N total"}\NormalTok{,}
    \StringTok{"Fallecimientos (eventos)"}\NormalTok{,}
    \StringTok{"Tiempo mediano de supervivencia (días)"}\NormalTok{,}
    \StringTok{"Edad media — x2 (años)"}\NormalTok{,}
    \StringTok{"Presión arterial sistólica media — x3 (mmHg)"}\NormalTok{,}
    \StringTok{"IMC medio — x5 (kg/m²)"}\NormalTok{,}
    \StringTok{"Tratamiento estándar — x1=1 (\%)"}\NormalTok{,}
    \StringTok{"Mujeres — x4=1 (\%)"}
\NormalTok{  ),}
  \AttributeTok{Valor =} \FunctionTok{c}\NormalTok{(}
    \FunctionTok{nrow}\NormalTok{(datos),}
    \FunctionTok{paste0}\NormalTok{(}\FunctionTok{sum}\NormalTok{(datos}\SpecialCharTok{$}\NormalTok{Status), }\StringTok{" ("}\NormalTok{, }\FunctionTok{round}\NormalTok{(}\FunctionTok{mean}\NormalTok{(datos}\SpecialCharTok{$}\NormalTok{Status) }\SpecialCharTok{*} \DecValTok{100}\NormalTok{, }\DecValTok{1}\NormalTok{), }\StringTok{"\%)"}\NormalTok{),}
    \FunctionTok{median}\NormalTok{(datos}\SpecialCharTok{$}\NormalTok{Survt),}
    \FunctionTok{round}\NormalTok{(}\FunctionTok{mean}\NormalTok{(datos}\SpecialCharTok{$}\NormalTok{x2), }\DecValTok{1}\NormalTok{),}
    \FunctionTok{round}\NormalTok{(}\FunctionTok{mean}\NormalTok{(datos}\SpecialCharTok{$}\NormalTok{x3), }\DecValTok{1}\NormalTok{),}
    \FunctionTok{round}\NormalTok{(}\FunctionTok{mean}\NormalTok{(datos}\SpecialCharTok{$}\NormalTok{x5), }\DecValTok{1}\NormalTok{),}
    \FunctionTok{paste0}\NormalTok{(}\FunctionTok{round}\NormalTok{(}\FunctionTok{mean}\NormalTok{(datos}\SpecialCharTok{$}\NormalTok{x1) }\SpecialCharTok{*} \DecValTok{100}\NormalTok{, }\DecValTok{1}\NormalTok{), }\StringTok{"\%"}\NormalTok{),}
    \FunctionTok{paste0}\NormalTok{(}\FunctionTok{round}\NormalTok{(}\FunctionTok{mean}\NormalTok{(datos}\SpecialCharTok{$}\NormalTok{x4) }\SpecialCharTok{*} \DecValTok{100}\NormalTok{, }\DecValTok{1}\NormalTok{), }\StringTok{"\%"}\NormalTok{)}
\NormalTok{  )}
\NormalTok{) }\SpecialCharTok{\%\textgreater{}\%}
  \FunctionTok{kable}\NormalTok{(}\AttributeTok{caption =} \StringTok{"Tabla 1. Características basales de la muestra"}\NormalTok{) }\SpecialCharTok{\%\textgreater{}\%}
  \FunctionTok{kable\_styling}\NormalTok{(}
    \AttributeTok{bootstrap\_options =} \FunctionTok{c}\NormalTok{(}\StringTok{"striped"}\NormalTok{, }\StringTok{"hover"}\NormalTok{, }\StringTok{"condensed"}\NormalTok{),}
    \AttributeTok{full\_width =} \ConstantTok{FALSE}
\NormalTok{  )}
\end{Highlighting}
\end{Shaded}

\begin{longtable}[t]{ll}
\caption{\label{tab:tabla-descriptiva}Tabla 1. Características basales de la muestra}\\
\toprule
Variable & Valor\\
\midrule
N total & 2000\\
Fallecimientos (eventos) & 1595 (79.8\%)\\
Tiempo mediano de supervivencia (días) & 370.068354106043\\
Edad media — x2 (años) & 5.5\\
Presión arterial sistólica media — x3 (mmHg) & 4\\
\addlinespace
IMC medio — x5 (kg/m²) & 5.5\\
Tratamiento estándar — x1=1 (\%) & 43.3\%\\
Mujeres — x4=1 (\%) & 151.3\%\\
\bottomrule
\end{longtable}

\hypertarget{curvas-de-kaplan-meier}{%
\subsection{Curvas de Kaplan-Meier}\label{curvas-de-kaplan-meier}}

\begin{Shaded}
\begin{Highlighting}[]
\CommentTok{\# Definición del objeto de supervivencia}
\NormalTok{surv\_obj }\OtherTok{\textless{}{-}} \FunctionTok{Surv}\NormalTok{(}\AttributeTok{time =}\NormalTok{ datos}\SpecialCharTok{$}\NormalTok{Survt, }\AttributeTok{event =}\NormalTok{ datos}\SpecialCharTok{$}\NormalTok{Status)}

\CommentTok{\# Estimación global}
\NormalTok{km\_global }\OtherTok{\textless{}{-}} \FunctionTok{survfit}\NormalTok{(surv\_obj }\SpecialCharTok{\textasciitilde{}} \DecValTok{1}\NormalTok{, }\AttributeTok{data =}\NormalTok{ datos)}

\CommentTok{\# Crear el gráfico y asignarlo a una variable}
\NormalTok{p\_km\_global }\OtherTok{\textless{}{-}} \FunctionTok{ggsurvplot}\NormalTok{(}
\NormalTok{  km\_global,}
  \AttributeTok{data =}\NormalTok{ datos,}
  \AttributeTok{conf.int =} \ConstantTok{TRUE}\NormalTok{,}
  \AttributeTok{risk.table =} \ConstantTok{TRUE}\NormalTok{,}
  \AttributeTok{xlab =} \StringTok{"Tiempo (días)"}\NormalTok{,}
  \AttributeTok{ylab =} \StringTok{"Probabilidad de supervivencia"}\NormalTok{,}
  \AttributeTok{title =} \StringTok{"Supervivencia global tras el ictus"}\NormalTok{,}
  \AttributeTok{ggtheme =} \FunctionTok{theme\_bw}\NormalTok{(),}
  \AttributeTok{palette =} \StringTok{"\#2E86AB"}\NormalTok{,}
  \AttributeTok{risk.table.height =} \FloatTok{0.25}
\NormalTok{)}

\CommentTok{\# Imprimir el gráfico}
\FunctionTok{print}\NormalTok{(p\_km\_global)}
\end{Highlighting}
\end{Shaded}

\begin{figure}

{\centering \includegraphics{trabajo_supervivencia_parte2_files/figure-latex/km-global-1} 

}

\caption{Figura 1. Curva de Kaplan-Meier global}\label{fig:km-global}
\end{figure}

\begin{Shaded}
\begin{Highlighting}[]
\NormalTok{km\_tto }\OtherTok{\textless{}{-}} \FunctionTok{survfit}\NormalTok{(surv\_obj }\SpecialCharTok{\textasciitilde{}}\NormalTok{ x1, }\AttributeTok{data =}\NormalTok{ datos)}

\CommentTok{\# Crear el gráfico y asignarlo a una variable}
\NormalTok{p\_km\_tto }\OtherTok{\textless{}{-}} \FunctionTok{ggsurvplot}\NormalTok{(}
\NormalTok{  km\_tto,}
  \AttributeTok{data =}\NormalTok{ datos,}
  \AttributeTok{conf.int =} \ConstantTok{TRUE}\NormalTok{,}
  \AttributeTok{risk.table =} \ConstantTok{TRUE}\NormalTok{,}
  \AttributeTok{pval =} \ConstantTok{TRUE}\NormalTok{,}
  \AttributeTok{pval.method =} \ConstantTok{TRUE}\NormalTok{,}
  \AttributeTok{xlab =} \StringTok{"Tiempo (días)"}\NormalTok{,}
  \AttributeTok{ylab =} \StringTok{"Probabilidad de supervivencia"}\NormalTok{,}
  \AttributeTok{title =} \StringTok{"Supervivencia por grupo de tratamiento"}\NormalTok{,}
  \AttributeTok{legend.labs =} \FunctionTok{c}\NormalTok{(}\StringTok{"Nuevo tratamiento (x1=0)"}\NormalTok{, }\StringTok{"Tratamiento estándar (x1=1)"}\NormalTok{),}
  \AttributeTok{legend.title =} \StringTok{"Tratamiento"}\NormalTok{,}
  \AttributeTok{ggtheme =} \FunctionTok{theme\_bw}\NormalTok{(),}
  \AttributeTok{palette =} \FunctionTok{c}\NormalTok{(}\StringTok{"\#2E86AB"}\NormalTok{, }\StringTok{"\#E84855"}\NormalTok{),}
  \AttributeTok{risk.table.height =} \FloatTok{0.28}
\NormalTok{)}

\CommentTok{\# Imprimir el gráfico}
\FunctionTok{print}\NormalTok{(p\_km\_tto)}
\end{Highlighting}
\end{Shaded}

\begin{figure}

{\centering \includegraphics{trabajo_supervivencia_parte2_files/figure-latex/km-tratamiento-1} 

}

\caption{Figura 2. Kaplan-Meier por grupo de tratamiento}\label{fig:km-tratamiento}
\end{figure}

\begin{center}\rule{0.5\linewidth}{0.5pt}\end{center}

\begin{quote}
\textbf{Nota:} Los apartados siguientes del trabajo desarrollarán el
ajuste e interpretación del modelo de Cox estándar, el modelo
estratificado y el modelo ampliado, una vez cargados y explorados los
datos reales del fichero \texttt{datos.txt}.
\end{quote}

\begin{center}\rule{0.5\linewidth}{0.5pt}\end{center}

\hypertarget{referencias}{%
\section{Referencias}\label{referencias}}

\begin{itemize}
\tightlist
\item
  Cox, D. R. (1972). Regression models and life-tables. \emph{Journal of
  the Royal Statistical Society: Series B}, 34(2), 187--202.
\item
  Therneau, T. M., \& Grambsch, P. M. (2000). \emph{Modeling Survival
  Data: Extending the Cox Model}. Springer.
\item
  Kassambara, A., Kosinski, M., \& Biecek, P. (2021). \emph{survminer:
  Drawing Survival Curves using `ggplot2'}. R package version 0.4.9.
\item
  R Core Team (2024). \emph{R: A language and environment for
  statistical computing}. R Foundation for Statistical Computing,
  Vienna, Austria.
\end{itemize}

\end{document}
